\section{System Design: Astrolabe}
\label{sec:system}


Astrolabe implements the categorical data model of Section~\ref{sec:data-model} as a desktop application. This section describes four architectural decisions that follow from the model's structure: the core/plugin boundary, the functor interpretation of plugins, the role of AI, and the component DSL for content.


\subsection{Minimal Core + Plugin Architecture}
\label{sec:sys-core}

The core system implements exactly three capabilities: (1)~obj/mor CRUD with a field whitelist and extensible metadata, (2)~Canvas rendering with sort-driven color and style mapping, and (3)~the sort system itself (automatic color assignment, core vocabulary, open extension). Everything else---analysis algorithms, data import, AI skills---is a plugin.

This boundary is not an arbitrary modularity choice; it reflects the categorical model. The core \emph{is} the Astrolabe category $\mathcal{A}(G)$: it maintains the objects, morphisms, sort labels, and identifiers that define the category. Plugins are operations \emph{on} the category---functors that bring external data in, enrichments that attach computed metadata, skills that help users construct new entities. The core knows nothing about PageRank, Lean, or community detection; it knows only about obj, mor, sort, and hash.

The practical consequence is domain agnosticism. Analysis methods that are meaningful for mathematical dependency networks (betweenness centrality, DAG depth) may be irrelevant for legal case hierarchies or biological pathway maps. Import formats differ by data source. If these capabilities were hardcoded, the system would be math-specific. The plugin boundary ensures that the same core serves any knowledge domain; plugins specialize it.

The split has a measurable effect on system complexity. The initial monolithic server (2515 lines, 63 API routes) was refactored into a 1008-line core (26 routes for CRUD, canvas state, and document access) plus plugin-registered routers (25 analysis endpoints, 1 import endpoint). New analysis algorithms or data sources require no changes to the core.


\subsection{Plugins as Functors}
\label{sec:sys-plugins-functors}

The correspondence between plugins and categorical constructions (Definition~\ref{def:plugin-functors}) is not merely an analogy---it is the design principle that determines the plugin interface.

Import plugins implement functors $F: \mathcal{E} \to \mathcal{A}(G)$. The ilean parser, our first import plugin, realizes $F_{\mathrm{Lean}}: \mathcal{L} \to \mathcal{A}(G)$. It maps Lean's \texttt{ConstantInfo} records to objects (mapping declaration kind to sort, sorry status to the \texttt{status} field), type dependencies to morphisms (with sort \texttt{uses}), and Lean qualified names to $\mathrm{Hex}_{12}$ identifiers via SHA-256 truncation. The functor property holds: if declaration~$A$ depends on~$B$ and $B$ depends on~$C$ in Lean, the composite dependency $A \to C$ is recoverable as a composed morphism in $\mathcal{A}(G)$. The output is standard obj/mor, subject to the same rendering, analysis, and cross-referencing as manually created entities.

Analysis plugins implement enrichments on $\mathcal{A}(G)$. They read the graph's structure, compute derived properties---centrality scores, community labels, bottleneck indicators---and write results back as metadata. They do not add or remove entities; they annotate existing ones. The analysis store accepts arbitrary key-value output from any plugin; the visual mapping layer reads from it dynamically, allowing a plugin's results to drive node size, node color, or edge style without any change to the rendering code.

The plugin interface is intentionally narrow: a plugin declares its name, version, endpoints, and skills in a JSON manifest; the loader scans, imports the entry module, and registers its API router. The narrow surface limits both the integration effort (a minimal plugin is ${\sim}30$ lines) and the potential for unintended side effects.


\subsection{AI as Proposal Generator}
\label{sec:sys-ai}

AI integration follows a strict \emph{proposal--confirmation} pattern. The user issues a natural language request or selects a built-in skill (e.g., \texttt{/add-node}). The AI generates a structured proposal---a JSON block containing obj/mor data with appropriate sort labels and content fields. The system detects these blocks in the AI's streaming output and renders them as interactive widgets. The user reviews the proposal and explicitly confirms it; only then is the data written to the knowledge graph.

This design reflects two principles. First, \emph{data integrity}: the knowledge graph is the user's intellectual artifact. Allowing unsupervised AI modification risks silent corruption---wrong sort assignments, spurious edges, hallucinated content. The proposal pattern ensures that every entity in the graph was either manually created or explicitly approved. Second, \emph{auditability}: there is no state in which an entity exists in the graph without the user's knowledge of how it got there.

The AI also serves as a code generator for the MDX component DSL. When a user requests ``show the proof tree for this theorem,'' the AI generates \texttt{<ProofTree root="abc123"/>}---exercising the DSL through natural language rather than requiring the user to know the component syntax.


\subsection{MDX Components as Knowledge Reference DSL}
\label{sec:sys-dsl}

The MDX content substrate supports a domain-specific language for referencing knowledge graph entities. Basic components provide live references: \texttt{<ObjBlock id="...">} renders an object as a card displaying its name, sort badge, statement, and proof; \texttt{<ObjRef id="...">} renders an inline reference with hover preview and click navigation. These components subscribe to the data store and reflect changes to the underlying object in real time---they are not static text snapshots.

Composite components, partially implemented, address higher-level queries: \texttt{<NodeCluster id="...">} renders an object with all its inclusion children as a collapsible tree; \texttt{<ProofTree root="...">} traces the \texttt{uses} dependency graph rooted at a theorem; \texttt{<FormalizationStatus id="...">} computes what fraction of a natural-language theorem's dependencies have corresponding \texttt{formalizes} morphisms.

The DSL framing clarifies the system's extensibility model. In this analogy, obj/mor are data types, sorts are the type system, MDX components are syntax, plugins are the standard library, and AI is the code generator. Adding a new composite component is adding a new syntactic form---it reads from the same underlying category and requires no changes to the data model or the core system.
